%% chapters/ch01-introduction.tex

\chapter{Introduction générale}
\label{ch:introduction}

\todo{Rédiger l'introduction générale.}

\section{Parcours et contexte}
\label{sec:intro:parcours}

Ce chapitre introductif présente le parcours scientifique et professionnel qui a conduit
à la rédaction de ce mémoire d'habilitation à diriger des recherches (HDR).

Depuis l'obtention de mon doctorat en \todo{année}, mes activités de recherche se sont
développées autour de la thématique centrale de \todo{thématique}.

\section{Organisation du manuscrit}
\label{sec:intro:organisation}

Ce manuscrit est organisé en \todo{N} parties principales~:

\begin{description}
  \item[Chapitre~\ref{ch:context}] présente le contexte scientifique et l'état de l'art
    dans lequel s'inscrivent mes travaux.
  \item[Chapitre~\ref{ch:contributions}] expose mes contributions scientifiques
    principales.
  \item[Chapitre~\ref{ch:synthesis}] propose une synthèse transversale de l'ensemble
    des travaux.
  \item[Chapitre~\ref{ch:perspectives}] trace les perspectives de recherche envisagées.
\end{description}

\section{Activités d'encadrement}
\label{sec:intro:encadrement}

Le tableau~\ref{tab:students} liste les doctorants et stagiaires que j'ai encadrés
depuis le début de ma carrière académique.

\begin{table}[htbp]
  \centering
  \caption{Étudiants encadrés}
  \label{tab:students}
  \begin{tabular}{llll}
    \toprule
    \textbf{Nom} & \textbf{Période} & \textbf{Niveau} & \textbf{Sujet} \\
    \midrule
    \studentrow{Étudiant A}{2018--2021}{Doctorat}{Sujet de thèse A}
    \studentrow{Étudiant B}{2020--2023}{Doctorat}{Sujet de thèse B}
    \studentrow{Étudiant C}{2022}{Master 2}{Sujet de stage}
    \bottomrule
  \end{tabular}
\end{table}
